\documentclass[12pt]{article}
\usepackage[utf8]{inputenc}
\usepackage[T1]{fontenc}
\usepackage{hyperref}
\title{Shina Language Resources for K-12 Curriculum Development}
\author{Compiled by Claude Sonnet 5.5 (claude-sonnet-5-5) for Andr\'as Kornai}
\date{}
\begin{document}
\maketitle

\section{Introduction}
Shina (ISO 639-3 scl; Glottocode shin1264) is a Dardic Indo-Aryan language of Gilgit-Baltistan and Kohistan in Pakistan, and of Dras and the Gurez valley in Indian-administered Kashmir. It has several dialects (Gilgiti, Kohistani, Guresi, Drasi). It is usually written in an extended Perso-Arabic (Nastaliq) script, and sometimes in Devanagari. Orthography was reported as standardized only in the late 2010s, with separate conventions in Gilgit-Baltistan/Kohistan and in Dras.

\section{Primary Educational Resources}
\begin{itemize}
\item Report on a Shina primer prepared for the Gilgit-Baltistan government and taught to nursery students from February 2017 (cited in the Wikipedia article). \url{https://en.wikipedia.org/wiki/Shina_language}
\item Pamir Times (2017) on a Gilgit-Baltistan plan to prepare syllabi for five regional languages including Shina as optional subjects, with a research academy. I found no evidence of its implementation. \url{https://pamirtimes.net/2017/09/14/gilgit-baltistans-5-languages-to-be-taught-at-schools-research-academy-to-be-set-up/}
\item Kashmir Despatch report that the J and K Board of School Education plans Shina textbooks for Classes I and II in the Gurez valley from the 2026 session. \url{https://kashmirdespatch.com/jkbose-move-rekindles-shinas-voice-in-gurez/}
\end{itemize}

\section{Linguistic and Reference Materials}
\begin{itemize}
\item Wikipedia, ``Shina language''. \url{https://en.wikipedia.org/wiki/Shina_language}
\item Wikipedia, ``Kohistani Shina''. \url{https://en.wikipedia.org/wiki/Kohistani_Shina}
\item T. Grahame Bailey, \emph{Grammar of the Shina language} (1924), as listed in the ASJP source record. \url{https://asjp.clld.org/sources/8277}
\item Pamir Times (2021) on Razwal Kohistani's three-volume Shina-Kohistani and Urdu dictionary, published by the Gandhara Hindku Academy. \url{https://pamirtimes.net/2021/10/22/dictionary-of-kohistani-shina-urdu-published/}
\item Omniglot page on Shina writing. \url{https://omniglot.com/writing/shina.htm}
\end{itemize}

\section{Religious and Community Resources}
No Shina scripture or specific community organization was identified in my search results.

\section{Digital and Audio Resources}
\begin{itemize}
\item Forum for Language Initiatives article on Shina language preservation efforts. \url{https://www.fli-online.org/documents/languages/shina/dying-shina.htm}
\end{itemize}
No audio or software resources were found in search results.

\section{Curriculum Development Context}
Shina is not firmly established in schools, though recent developments are notable: a 2017 nursery primer in Gilgit-Baltistan, a 2017 plan for optional regional-language syllabi, and planned Class I and II textbooks in Gurez from 2026. A commentary in Business Recorder argues that Gilgit-Baltistan has still not included regional languages in school curricula. \url{https://www.brecorder.com/news/40272925/gilgit-baltistan-preserve-the-linguistic-diversity} Dialect and script divergence between regions is the main design issue. The dictionary and Bailey's grammar are the most useful references.
\end{document}